% !TEX root = chapter-standalone.tex

\devel{
\section{Subcategories}
\label{sec:subcategories}

\todographicsjira{429}{\alphubel: @Gioele: Needs graphics for subcategorie}

In \cref{chap:new-from-old-algebra} we saw that a \SY{semigroup}, \SY{monoid}, or \SY{group} can contain smaller structures of the same kind: \SY{subsemigroups}, \SY{submonoids}, and \SY{subgroups}.
The same phenomenon occurs for categories.
Frequently, a selection of some of the objects and some of the morphisms of a category~\CatC is ``self-sufficient'': composing selected morphisms never leads outside the selection, and the selected objects bring their \SY{identity morphisms} with them.
Such a selection then forms a category in its own right, called a \emph{subcategory} of~\CatC.

As with sub-things in algebra, subcategories are useful for modeling \emph{restrictions}.
Given a category describing a family of systems and processes, we often want to work with a restricted class of processes---only those which are reversible, or feasible, or within budget, or compliant with a specification.
Whenever the restricted class is closed under composition and contains identities, the restriction is again a category, and everything we know about categories applies to it.

\begin{ctdefinition}[Subcategory]
    \label{def:subcategory}
    \SYNDEF{subcategory}
    Let~\CatC be a \SY{semicategory}.
    A \emph{subsemicategory}~\CatD of~\CatC is specified by:

    \constit

    \begin{enumerate}
        \item A sub-collection~$\ObD$ of the objects of~$\CatC$;
        \item For any two objects~$\Obja,\Objb\setin \ObD$, a subset of the morphisms of \CatC between them:
              \begin{equation}
                  \HomSet{\CatD}{\Obja}{\Objb}\setsubseteq \HomSet{\CatC}{ \Obja}{\Objb}.
              \end{equation}
    \end{enumerate}

    \condit

    \begin{enumerate}
        \item Closure under composition: if~$\mora \colon \Obja \mto \Objb$ and~$\morb\colon \Objb\mto \Objc$ are morphisms of \CatD, then the composite~$\morab$, computed in~\CatC, is again a morphism of~\CatD.
    \end{enumerate}

    If~\CatC is a category, a \emph{subcategory}~\CatD of~\CatC is a subsemicategory which additionally satisfies:

    \begin{enumerate}
        \setcounter{enumi}{1}
        \item Identities: if~$\Obja\setin \ObD$, then the \SY{identity morphism} $\catidat\Obja \setin \HomSet{\CatC}{\Obja}{\Obja}$ is also in~$\HomSet{\CatD}{\Obja}{\Obja}$.
    \end{enumerate}
\end{ctdefinition}

\begin{remark}
    A subsemicategory~\CatD of a \SY{semicategory}~\CatC is itself a \SY{semicategory}: composition restricts to the selected morphisms by the closure condition, and associativity is inherited from~\CatC, since it holds for \emph{all} triples of composable morphisms of~\CatC, in particular the selected ones.
    Likewise, a \SY{subcategory} of a category is itself a category, with the inherited \SY{identity morphisms}.
    Compare this with the corresponding observation for \SY{subsemigroups} (\cref{rem:subsemigroup-is-semigroup}).
\end{remark}

There are two extreme ways of building subcategories, which both occur so often that they have names: we can keep all the objects and restrict only the morphisms, or restrict the objects and keep all morphisms between the ones we kept.

\begin{ctdefinition}[Wide and full subcategories]
    \label{def:wide-full-subcategory}
    Let~\CatD be a \SY{subcategory} of a category~\CatC.
    \begin{enumerate}
        \item \SYNDEF{wide subcategory}\CatD is called \emph{wide} if it contains all objects of~\CatC, that is, $\ObD = \ObC$.
        \item \SYNDEF{full subcategory}\CatD is called \emph{full} if, for all~$\Obja, \Objb \setin \ObD$, it holds that $\HomSet{\CatD}{\Obja}{\Objb} = \HomSet{\CatC}{\Obja}{\Objb}$.
    \end{enumerate}
\end{ctdefinition}

\begin{remark}
    A \SYN{full subcategory}{full} subcategory is completely determined by the choice of which objects to keep: the morphisms simply come along.
    Choosing a full subcategory therefore amounts to focusing attention on a class of objects of interest, without restricting the processes between them.
    By contrast, choosing a wide subcategory amounts to restricting the admissible processes, while keeping all objects in play.
\end{remark}

\subsection{Subcategories of \Set and \Rel}

\begin{example}[\FinSet]
    \SYNDEF{category of finite sets and functions}
    \label{ex:FinSet}
    \FinSet is the category of finite sets and all functions between them.
    It is a \SY{subcategory} of the category \Set of sets and functions.
    While an object~$\Obja \setin \Obof\Set$ is a set with arbitrary cardinality,~$\Obof{\FinSet}$ only includes sets which have finitely many elements.
    Objects of \FinSet are in \Set, but the converse is not true.
    Furthermore, given~$\Obja,\Objb\setin \Obof\FinSet$, we take~$\HomSet{\FinSet}{\Obja}{\Objb} =\HomSet{\Set}{\Obja}{\Objb}$.
    In particular, \FinSet is a \SYN{full subcategory}{full} subcategory of \Set (but not a wide one).
\end{example}

\begin{example}[\Set and \Rel]
    The category \Set is a \SY{subcategory} of \Rel.
    To show this, we need to prove the conditions presented in \cref{def:subcategory}.
    \begin{enumerate}
        \item In both \Rel and \Set, the collection of objects is all sets.
        \item Given~$\Obja,\Objb\setin \Obof{\Set}$, we know that~$\HomSet{\Set}{\Obja}{\Objb}\setsubseteq \HomSet{\Rel}{\Obja}{\Objb}$, since all functions between sets~$\Obja,\Objb$ are a particular subset of all relations between~$\Obja,\Objb$.
        \item Let~$\relA\setsubseteq \Obja\cartprod \Objb$ and~$\relB\setsubseteq \Objb \cartprod \Objc$ be relations which are functions.
              We need to show that their composition in \Rel, expressed as~$\relA\mthen \relB\setsubseteq \Obja\cartprod \Objc$, is again a function.
              This was proven in \cref{lem:comprelfun}.
        \item For each~$\Obja \setin \Obof{\Set}$, the identity relation~$\catidat\Obja=\makeset{\tupp{\ela,\ela\elprime}\setin {\Obja} \cartprod \Obja \mid \ela=\ela\elprime}$ corresponds to the identity function~$\catidat\Obja \colon \Obja \mto \Obja$ in \Set.
    \end{enumerate}
    Note that \Set is a \SYN{wide subcategory}{wide} subcategory of \Rel: both have all sets as objects.
    It is not a full one: not every relation is a function.
\end{example}

\begin{example}
    \label{exa:bijections-wide-subcategory}
    Keep all sets as objects, but keep only the \SY{bijective} functions as morphisms.
    This is a \SYN{wide subcategory}{wide} subcategory of \Set: the composition of two bijections is a bijection, and identity functions are bijections.

    This is the pattern behind many subcategories that arise in practice: one restricts to the morphisms that are \emph{reversible}.
    For an applied reading, if objects are file formats and morphisms are conversion routines, the restriction to bijections corresponds to considering only \emph{lossless} conversions.
\end{example}

\subsection{Non-examples and near-misses}

Being a \SY{subcategory} is a genuine condition on a selection of objects and morphisms; it is instructive to see how a selection can fail.

\begin{example}[Failure of closure under composition]
    \label{exa:quadratics-closure-fails}
    Consider the selection with a single object, the set~\reals, and with morphisms the polynomial functions~$\reals \sto \reals$ of degree at most~$2$.
    This selection is not closed under composition: composing the function~$\ela \mapsto \ela^2$ with itself yields~$\ela \mapsto \ela^4$, which has degree~$4$.
    So ``sets and low-degree polynomial maps'' is not a \SY{subcategory} of \Set.

    This failure mode is worth remembering in modeling practice: a class of ``simple'' components (here: quadratic response curves) is often \emph{not} stable under composition, and the composite system can be strictly more complex than its parts.
\end{example}

\begin{example}[Failure of the identity condition]
    \label{exa:noninjective-subsemicategory}
    In \Set, keep all sets, and keep as morphisms only the functions that are \emph{not} \SY{injective}.
    If~$\mora$ is not injective, then any composite~$\morab$ is also not injective, so this selection \emph{is} closed under composition, and it forms a subsemicategory of \Set.
    However, it is not a \SY{subcategory}: identity functions are injective, so no object has its \SY{identity morphism} in the selection.

    This is the categorical analogue of a situation we met in algebra: a subset of a \SY{monoid} may be closed under composition and yet fail to contain the \SY{neutral element} (\cref{exa:zero-not-submonoid}).
\end{example}

\subsection{Subcategories, sub-monoids, free categories}

\begin{example}
    \label{exa:subcats-of-monoid-cat}
    Recall from \cref{def:monoidcat} that any \SY{monoid}~$\monoidAdefinition$ can be viewed as a category~$\monoidcat{\monoidA}$ with a single object, whose morphisms are the elements of~$\monoidAset$.
    The \SY{subcategories} of~$\monoidcat{\monoidA}$ that contain the single object correspond exactly to the \SY{submonoids} of~$\monoidA$: a selection of morphisms closed under composition and containing the identity is nothing but a subset of~$\monoidAset$ closed under~$\mtimes$ and containing~$\idmon$.
    Likewise, subsemicategories correspond to \SY{subsemigroups}.

    For instance, viewing~$\tup{\wnumbers, +, 0}$ as a one-object category, the even numbers form a \SY{subcategory}, while the odd numbers fail on both counts: they are not closed under composition, and they do not include the identity~$0$.
\end{example}

\begin{example}
    \label{exa:subgraph-subcategory}
    Let~$\graphA$ be a graph, and let~$\graphA'$ be a subgraph (a selection of some vertices and some edges among them).
    Then the free category~$\pathcat{\graphA'}$ is a \SY{subcategory} of the free category~$\pathcat{\graphA}$ (\cref{sec:catsfromgraphs}): every path using only edges of~$\graphA'$ is in particular a path in~$\graphA$, concatenation is computed the same way in both, and trivial paths at vertices of~$\graphA'$ qualify.

    This kind of subcategory has a natural applied reading.
    If~$\graphA$ is the graph of an intermodal mobility network---vertices are locations, edges are individual road segments, subway rides, and so on---then morphisms of~$\pathcat{\graphA}$ are the possible trips.
    Selecting the subgraph~$\graphA'$ of walking edges only, the subcategory~$\pathcat{\graphA'}$ consists of the trips one can make on foot: a restriction of the mobility model obtained by disabling some modes of transport.
\end{example}

\begin{exercise}
    \label{ex:wide-full-classify}
    Classify the following \SY{subcategories} as \SYN{wide subcategory}{wide}, \SYN{full subcategory}{full}, both, or neither:
    \begin{enumerate}
        \item \FinSet inside \Set;
        \item \Set inside \Rel;
        \item the subcategory of sets and bijections (\cref{exa:bijections-wide-subcategory}) inside \Set;
        \item \CatC inside \CatC (any category as a subcategory of itself).
    \end{enumerate}
\end{exercise}
\begin{solution}
    (1) Full, not wide: only finite sets are kept, but all functions between them.
    (2) Wide, not full: all sets are kept, but only those relations which are functions.
    (3) Wide, not full: all sets are kept, but (in general) not all functions.
    (4) Both wide and full.
\end{solution}




\devel{

\section{Further materials parked here for now}

\subsection{Subcategories of endomorphisms}
\label{sec:Draw}

\todographics{@Gioele: put here a figure from the exercises to illustrate the concept of objects and morphisms in Draw}
\begin{definition}[Drawings]
    \label{def:Draw}
    \SYNDEF{category of drawings}
    There exists a category~\Draw in which:
    \begin{enumerate}
        \item An object in~$\alpha\setin \Obof\Draw$ is a black-and-white drawing,
              that is a function~$\alpha \colon \reals^2 \sto \boolset$.
        \item A morphism in~$\HomSet\Draw{\alpha}{\beta}$ between two drawings~$\alpha$ and~$\beta$ is an invertible map~$f\colon \reals^2 \sto \reals^2$ such that~$\alpha(x) = \beta(f(x))$.
        \item The identity function at any object~$\alpha$ is the identity map on~$\reals^2$.
        \item Composition is given by function composition.
    \end{enumerate}
\end{definition}

\begin{exercise}
    \label{ex:draw}
    Check whether just considering
    \begin{itemize}
        \item affine invertible transformations, or
        \item rototranslations, or
        \item scalings, or
        \item translations, or
        \item rotations,
    \end{itemize}
    as morphisms forms a \SY{subcategory} of~\Draw.
\end{exercise}
\begin{solution}
    \begin{marginfigure}
        \begin{center}
            \includesag{scaling}
        \end{center}
        \caption{Example of scaling.}
    \end{marginfigure}

    \begin{marginfigure}
        \begin{center}
            \includesag{translation}
        \end{center}
        \caption{Example of translation.}
    \end{marginfigure}

    \begin{marginfigure}
        \begin{center}
            \includesag{rotation}
        \end{center}
        \caption{Example of rotation.}
    \end{marginfigure}

    \begin{marginfigure}
        \begin{center}
            \includesag{rototranslation}
        \end{center}
        \caption{Example of rototranslation.}
    \end{marginfigure}

    \begin{marginfigure}
        \begin{center}
            \includesag{affinetrafo}
        \end{center}
        \caption{Example of affine transformation with~$A=1.5\begin{bmatrix}
                    \cos(\pi/4) & \sin(\pi/4) \\-\sin(\pi/4)& \cos(\pi/4)
                \end{bmatrix}$.}
    \end{marginfigure}
    We check the specializations one by one.
    In all specializations, we consider the same objects as in~\Draw.
    \begin{itemize}
        \item \textbf{Scalings.}
              Let~$s,t\setin \reals$.
              Scalings can be represented as functions of the form
              \begin{equation}
                  \defmapperiod{
                      \scaling_{s,t}
                  }{
                      \reals^2
                  }{
                      \to
                  }{
                      \reals^2
                  }{
                      \tup{x,y}
                  }{
                      \tup{s x,t y}
                  }
              \end{equation}
              By just considering morphisms which are scalings, we are considering a subset of all morphisms.
              Furthermore, the composition of two scalings is again a scaling.
              Indeed, consider scalings~$\scaling_{s,t}$,~$\scaling_{u,v}$.
              We have
              \begin{equation}
                  \begin{aligned}
                      (\scaling_{s,t}\mthen \scaling_{u,v})(x,y)
                       & =\scaling_{u,v}(sx, ty) \\
                       & =\tup{usx, vty} \\
                       & =\scaling_{us,vt}.
                  \end{aligned}
              \end{equation}
              Finally, the \SY{identity morphism} in \Draw corresponds to a scaling of the form $\scaling_{1,1}$.
        \item \textbf{Translations.}
              Let~$s,t\setin \reals$.
              Translations are functions of the form
              \begin{equation}
                  \defmapperiod{
                      \translation_{s,t}
                  }{
                      \reals^2
                  }{
                      \to
                  }{
                      \reals^2
                  }{
                      \tup{x,y}
                  }{
                      \tup{x+s,y+t}
                  }
              \end{equation}
              By just considering morphisms which are translations, we are considering a subset of all morphisms.
              Furthermore, the composition of two translations is again a translation.
              Indeed, consider scalings~$\translation_{s,t}$,~$\translation_{u,v}$.
              We have
              \begin{equation}
                  \begin{aligned}
                      (\translation_{s,t}\mthen \translation_{u,v})(x,y)
                       & =\translation_{u,v}(x+s, y+t) \\
                       & =\tup{x+s+u, y+t+v} \\
                       & =\translation_{s+u,t+v}.
                  \end{aligned}
              \end{equation}
              Finally, the \SY{identity morphism} in \Draw corresponds to a translation of the form $\translation_{0,0}$.
        \item \textbf{Rotations.}
              Let~$\theta \setin [0,2\pi)$.
              Rotations are functions of the form
              \begin{equation}
                  \defmapperiod{
                      \rotation_{\theta}
                  }{
                      \reals^2
                  }{
                      \to
                  }{
                      \reals^2
                  }{
                      \tup{x,y}
                  }{
                      \tup{x\cos(\theta)+y\sin(\theta), y\cos(\theta)-x\sin(\theta)}
                  }
              \end{equation}
              By just considering morphisms which are rotations, we are considering a subset of all morphisms.
              Furthermore, the composition of two rotations is again a rotation.
              Indeed, consider rotations~$\rotation_{\theta}$,~$\rotation_{\phi}$.
              We have
              \begin{equation}
                  \begin{aligned}
                      (\rotation_{\theta}\mthen \rotation_{\phi})(x,y)
                       & =\rotation(\tup{x\cos(\theta)+y\sin(\theta), y\cos(\theta)-x\sin(\theta)}) \\
                       & =\tup{x\cos(\theta+\phi)+y\sin(\theta+\phi), y\cos(\theta+\phi)-x\sin(\theta+\phi)} \\
                       & =\rotation_{\theta+\phi}.
                  \end{aligned}
              \end{equation}
              Finally, the \SY{identity morphism} in \Draw corresponds to a rotation of the form $\rotation_{0}$.
        \item \textbf{Rototranslations.}
              Let~$s,t\setin \reals$ and~$\theta \setin [0,2\pi)$.
              Rototranslations are functions arising from the combination of rotations and translations, and are therefore of the form:
              \begin{equation}
                  \defmapperiod{
                      \rototrans_{\theta,s,t}
                  }{
                      \reals^2
                  }{
                      \to
                  }{
                      \reals^2
                  }{
                      \tup{x,y}
                  }{
                      \tup{x\cos(\theta)+y\sin(\theta)+s, y\cos(\theta)-x\sin(\theta)+t}
                  }
              \end{equation}
              By just considering morphisms which are rotations, we are considering a subset of all morphisms.
              Furthermore, the composition of two rototranslations is again a rototranslation.
              Consider rototranslations~$\rototrans_{\theta,s,t}$,~$\rototrans_{\phi,u,v}$.
              We have:
              \begin{equation}
                  \begin{aligned}
                       & (\rototrans_{\theta,s,t}\mthen \rototrans_{\phi,u,v})(x,y) \\
                       & =\rototrans_{\phi,u,v}(x\cos(\theta)+y\sin(\theta)+s, y\cos(\theta)-x\sin(\theta)+t) \\
                       & =\langle(x\cos(\theta)+y\sin(\theta)+s)\cos(\phi)+(y\cos(\theta)-x\sin(\theta)+t))\sin(\phi)+u, \\
                       & (y\cos(\theta)-x\sin(\theta)+t)\cos(\phi) - (x\cos(\theta)+y\sin(\theta)+s)\sin(\phi)+v\rangle \\
                       & =\langle x\cos(\theta+\phi)+y\sin(\theta+\phi)+s\cos(\phi)+t\sin(\phi)+u, \\
                       & y\cos(\theta+\phi)-x\sin(\theta+\phi)+t\cos(\phi)-s\sin(\phi)+v\rangle \\
                       & =\rototrans_{\theta+\phi,s\cos(\phi)+t\sin(\phi)+u, t\cos(\phi)-s\sin(\phi)+v}(x,y).
                  \end{aligned}
              \end{equation}
              Finally, the \SY{identity morphism} in \Draw corresponds to a rotation of the form $\rototrans_{0,0,0}$.
        \item \textbf{Affine transformations.}
              Let~$\mat{A}\setin \reals^{\matdim{2}{2}}$ and~$\mat{b}\setin \reals^{\matdim{2}{1}}$.
              Affine transformations are functions that could arise from the combination of rotations and translations, and scalings, and are therefore of the form:
              \begin{equation}
                  \defmapcomma{
                  \affinetrafo_{\mat{A},\mat{b}}
                  }{
                  \reals^2
                  }{
                  \to
                  }{
                  \reals^2
                  }{
                  \tup{x,y}
                  }{
                  \tup{a_{11}x+a_{12}y+b_{11},a_{21}x+a_{22}y+b_{21}}
                  }
              \end{equation}
              where~$*_{ij}$ represents the element at the~$i$-th row and~$j$-th column of~$*$.

              Some special cases are:
              \begin{itemize}
                  \item With~$\mat{A}=\mathbb{1}$ we obtain translations~$\translation_{b_{11},b_{21}}$;
                  \item With~$\mat{b}=\begin{bmatrix}
                                0 & 0
                            \end{bmatrix}\mattransp$ and~$A=\begin{bmatrix}
                                s & 0 \\0& t
                            \end{bmatrix}$ we obtain scalings~$\scaling_{s,t}$;
                  \item With~$\mat{b}=\begin{bmatrix}
                                0 & 0
                            \end{bmatrix}\mattransp$ and~$A=\begin{bmatrix}
                                \cos(\theta) & \sin(\theta) \\-\sin(\theta)& \cos(\theta)
                            \end{bmatrix}$ we obtain rotations~$\rotation_{\theta}$.
              \end{itemize}
              By just considering morphisms which are affine transformations, we are considering a subset of all morphisms.
              Furthermore, the composition of two affine transformations is again an affine transformation.
              Clearly, the composition of affine transformations~$\affinetrafo_{\mat{A},\mat{b}}$,~$\affinetrafo_{\mat{C},\mat{d}}$ is~$\affinetrafo_{\mat{CA},\mat{Cb}+\mat{d}}$
              Finally, the \SY{identity morphism} in \Draw corresponds to an affine transformation of the form~$\affinetrafo_{\mathbb{1},\zeromat^{\matdim{2}{1}}}$.
    \end{itemize}
\end{solution}


\subsection{Subcategories of Berg}
\label{sec:subcat_berg}

Recall the category \Berg presented in \cref{sec:trekking}.
In the following, we want to give both a positive and a negative example of \SY{subcategories} related to \Berg.

%We first start our discussion by introducing an \emph{amateur} version of \Berg, called \Bergama, which only considers paths (morphisms) in \Berg, whose steepness does not exceed a critical value, say 1/2.
%Is \Bergama a subcategory of \Berg?
We start our discussion by introducing a \emph{limited} version of \Berg, called~$\Berg_\alpha$, which only considers paths (morphisms) in \Berg, whose steepness does not exceed the critical value~$\alpha\setin [0,1]$.
Is~$\Berg_\alpha$ a \SY{subcategory} of \Berg?
We check the different conditions:
\begin{enumerate}
    \item The constraint on the maximum steepness restricts the objects which are acceptable in~$\Berg_\alpha$ via the \SY{identity morphisms} of \Berg.
          Indeed, recall that given an object~$\tup{\styleobj{p},\styleobj{v}}\setin \Obof{\Berg}$, the \SY{identity morphism} is defined as~$\stylemorph{1}_{\tup{\styleobj{p},\styleobj{v}}}=\tup{\gamma,0}$, with~$\gamma(0)=p$ and~$\dot{\gamma}(0)=v$.
          The steepness is computed via~$\styleobj{v}$.
          In particular,~$\Berg_\alpha$ will only contain objects whose \SY{identity morphisms} do not exceed the steepness constraint, In other words~$\Obof{\Berg_\alpha} \setsubseteq \Obof\Berg$.
    \item For~$\Obja,\Objb\setin \Obof{\Bergama}$, we know that paths satisfying the steepness constraint are specific paths in \Berg:~$\Hom_{\Berg_\alpha}\setsubseteq \Hom_{\Berg}$.
    \item Given two morphisms~$\mora,\morb$ which can be composed in~$\Berg_\alpha$, the maximum steepness of their composition~$\morab$ is given by:
          \begin{equation}
              \MaxSteepness(\morab)
              =
              \max \makeset{
                  \MaxSteepness(\mora),
                  \MaxSteepness(\morb)
              }
              <
              \alpha.
          \end{equation}
    \item The  \SY{identity morphisms} in \Berg which satisfy the steepness constraint are, by definition, in~$\Berg_\alpha$ and they act as identities there.
\end{enumerate}

This shows that~$\Berg_\alpha$ is a \SY{subcategory} of \Berg.

What would an example of non-subcategory of \Berg be?
We could try defining a new category \Berglazy, which now discriminates morphisms based on the lengths of the paths they represent.
For instance, assume that as amateur hikers, we don't want to consider morphisms which are more than \unit[1]{km} long.
By concatenating two paths (morphisms) of length \unit[0.6]{km} in \Berglazy, the resulting composition will be \unit[1.2]{km}, violating the posed constraint and hence not being in \Berglazy.
This violates the third property of \cref{def:subcategory}.
}


\devel{
\subsection[Other examples]{Other examples of \SY{subcategories} in engineering}

In engineering, it is very common to look at specific types of functions; in many cases, the properties of a certain type of function are preserved by function composition, and so they form a category.

\subsubsection{\Injset is a \SY{subcategory} of \Set}
% \todostructure{\alphubel: We have alrady defined informally \SY{injective} function in the preliminaries.
%     Move this there?}
% \begin{definition}[Injective function]
%     \label{def:injective-function}
%     Let~$\mapa \colon \setA \sto \setB$ be a function.
%     The function~$\mapa$ is \emph{injective} if, for all~$\ela,\elb\setin \setA$ holds:
%     \begin{equation}
%         \prfperiod{
%             \mapa(\ela)=\mapa(\elb)
%         }{
%             \ela=\elb
%         }
%     \end{equation}
% \end{definition}

\begin{example}
    \SYNDEF{category of sets and injective functions}
    \label{ex:Injset}
    We can define a category \Injset that has the same objects as \Set but restricts the morphisms to be \emph{\SY{injective functions}} (\cref{def:injective-function}).
    We want to show that \Injset is a \SY{subcategory} of \Set.
    Composition and \SY{identity morphisms} are defined as in \Set.

    Since~$\Obof{\Injset}=\Obof{\Set}$, the first condition of \cref{def:subcategory} is satisfied.
    Injective functions are a particular type of functions: this satisfies the second condition.
    Given~$\Obja\setin \Obof{\Injset}$, the \SY{identity morphism}~$\catidat\Obja\setin \HomSet{\Set}{\Obja}{\Obja}$ corresponds to the \SY{identity morphism} in~$\HomSet{\Injset}{\Obja}{\Obja}$: the identity function is \SY{injective}.
    This proves the fourth condition.
    To check the third condition, consider two morphisms~$\mora \setin \HomSet{\Set}{\Obja}{\Objb}$,~$\morb \setin \HomSet{\Set}{\Objb}{\Objc}$ such that~$\mora \setin \HomSet{\Injset}{\Obja}{\Objb}$ and~$\morb\setin \HomSet{\Injset}{\Objb}{\Objc}$.
    From the \SY{injectivity} of~$\mora,\morb$, we know that given~$\ela,\ela\elprime\setin \Obja$,
    \begin{equation}
        \prfdoublecomma{
            \mora(\ela)=\mora(\ela\elprime)
        }{
            \ela=\ela\elprime
        }
    \end{equation}
    and~$\elb,\elb\elprime\setin \Obja$,
    \begin{equation}
        \prfdoubleperiod{
            \morb(\elb)=\morb(\elb\elprime)
        }{
            \elb=\elb\elprime
        }
    \end{equation}
    Furthermore, we have:
    \begin{equation}
        \prfdoublecomma{
            (\morab)(\ela)
            =(\morab)(\ela\elprime)
        }{
            \prftree[double]{
                \mora(\ela)=\mora(\ela\elprime)
            }{
                \ela=\ela\elprime
            }
        }
    \end{equation}
    which proves the third condition of \cref{def:subcategory}: the composition of \SY{injective} functions is \SY{injective}.
\end{example}

\devel{
    \todojira{462}{Put the following in monoids}
    \begin{definition}[Continuous functions]\label{def:continuous-function}

        Let~$\mapa \colon \reals\sto \reals$ be a function.
        We call~$\mapa$ \emph{continuous} at~$c\setin \reals$ if
        %
        \begin{equation}
            \lim_{x\to c}\mapa(x)=f(c);
        \end{equation}
        %
        $\mapa$ is continuous over~\reals if the condition is satisfied for all~$c\setin \reals$.

    \end{definition}

    \begin{example}
        We can define a category~\Cont in which
        \begin{equation}
            \Obof\Cont=\makeset{\reals, \reals^2, \reals^3, \dots }
        \end{equation}
        and in which the morphisms are given by continuous functions.
        Composition and identity are as in~\Set.
        We want to show that \Cont is a \SY{subcategory} of~\Set.
        \todojira{126}{write down formally and use that composition of continuous is continuous}
    \end{example}

    \todojira{127}{Differentiable functions: Set to Manifolds}

    \begin{definition}[Differentiable functions]\label{def:differentiable}
        A function $f\colon U\subset \reals\sto \reals$, defined on an open set $U$, is \emph{differentiable} at $a\setin U$ if the derivative
        \begin{equation}
            f'(a)=\lim_{h\to 0} \frac{f(a+h)-f(a)}{h}
        \end{equation}
        exists; $f$ is differentiable on $U$ if it is differentiable at every point of $U$.
    \end{definition}

    \begin{example}
        The composition of differentiable functions is differentiable.
    \end{example}

    \todojira{128}{Lipschitz bounded}

    \begin{definition}[Lipschitz continuous function]\label{def:Lipschitz}
        A real valued function $f\colon \reals\sto \reals$ is called \emph{Lipschitz} continuous if there exists a positive real constant $\kappa$ such that, for all $x_1,x_2\setin \reals$:
        \begin{equation}
            \absvalueof{f(x_1)-f(x_2)} \leq \kappa \absvalueof{ x_1-x_2}.
        \end{equation}
    \end{definition}

    \begin{example}
        The composition of differentiable functions is differentiable
    \end{example}

    \todojira{129}{smooth; cont diff, composition is compdiff}
    %\begin{exercise}
    %Check that \Set, as specified above, does in fact define a category.
    %\end{exercise}

    \subsubsection{Generalization outside \reals}
    Generalization to more general spaces.
    We used the fact that R is:
    \begin{itemize}
        \item For defining continuous functions, we used the fact that~\reals is topological space (minimum needed for defining a continuous function).
              In fact, the real definition of continuous function is:

              \todojira{130}{ADD: continuous function}

        \item To define Lipschitz we needed the fact that~ \reals is a metric space
        \item For differentiable, smooth, you define this on Manifolds.
              Exists tangent space.
    \end{itemize}

    Hence, in general, the objects of these are different, so it's not really a relation of \SY{subcategory}, that requires the objects to be the same.
    However, we can generalize this notion using \SY{functors}.

    \todojira{131}{functor F:C→D that is both \SY{injective} on objects and a \SY{faithful functor}.}
}

\todojira{461}{\alphubel: Add set and pos example for subcats.}
}



}
